\section{Related work}
\label{sec:related_target_centering}

Physical knowledge enters learned models through several interfaces. Hybrid
process models combine mechanistic structure with learned quantities
\citep{PsichogiosUngar1992}; physics-informed neural networks use
differential-equation residuals in the objective
\citep{RaissiPerdikarisKarniadakis2019}; APHYNITY learns a constrained
correction to incomplete dynamics \citep{YinEtAl2021}. These approaches motivate
studying which part of a prediction problem is assigned to the physical
description and which part remains statistical. Here the physical prescription
defines the response coordinate, and the flow learns its conditional density.

Residual targets have a direct precedent in $\Delta$-machine learning, where
a learned correction augments an inexpensive quantum-chemical baseline
\citep{RamakrishnanEtAl2015}. Conditional transformation models extend response
transformation to distributional regression by estimating transformations
that depend on explanatory variables \citep{HothornKneibBuhlmann2014}.
These lines connect baseline correction with conditional distribution
modeling. Our center is prescribed before flow fitting, so the experimental
variable is the residual task presented to a matched density estimator.

Response parameterization can affect both model adequacy and optimization.
Box and Cox studied transformations that make a simple regression and error
model appropriate, comparing likelihoods on the original response scale with
the Jacobian included \citep{BoxCox1964}. PopArt changes target normalization
while compensating the output layer to preserve the unnormalized predictor,
relating normalization to modified squared-loss SGD updates
\citep{VanHasseltEtAl2016}. For squared-loss kernel regression, target power
across kernel eigenfunctions links target--model alignment to generalization
\citep{CanatarBordelonPehlevan2021}. Together these results show why an
information-preserving coordinate choice can still matter to a restricted
learner. Our arms share a global target scale and fit a residual offset,
focusing the comparison on input-dependent centering.

Normalizing flows permit tractable density evaluation through change of
variables. Conditional masked autoregressive flows incorporate explanatory
inputs into their transformations \citep{PapamakariosPavlakouMurray2017},
and neural spline flows provide monotone rational-quadratic maps with analytic
inverses \citep{DurkanEtAl2019}. Flow expressiveness also depends on depth,
parameterization and numerical conditioning \citep{KoehlerMehtaRisteski2021}.
Numerical conditioning concerns the invertible map and its Jacobian, distinct
from the physical information supplied to the neural conditioner.

These studies establish residual modeling and coordinate choice as relevant
tools; they do not determine the ordering of finite spline-flow NLLs under
restricted neural inputs. Our contribution is the controlled information-axis
comparison for a prescribed physical target center. Matched raw-target
baselines separate centering benefit from the absolute benefit of additional
neural information, and regional evaluations test how aggregate changes are
distributed across the physical domain.
